% Distributivity shortens the unchanged shifted-affine certificate.
\section{Factoring the packed mask: 93 operations}
\label{sec:factored-mask}

The system and fixed indices of Section~\ref{sec:shifted-affine}
admit a shorter arithmetic certificate by factoring the indicator
code out of two terms of the packed mask. No source equation,
encoding, admissibility condition, or positive unknown changes.

\begin{theorem}\label{thm:best93}
The universal system of Theorem~\ref{thm:best94}, with exactly
the same admissible fixed indices and positive witnesses, admits
a straight-line certificate of 93 operations: 50 multiplications
and 43 additions. There remain 34 positive unknowns and 22
equations. Its supplied parameters are $(x,V,H,T_L)$, with
three fixed index parameters besides the positive input $x$.
The literal numerals are $\{1,3,8,15\}$; fixed numerals and
equality tests are free.
\end{theorem}

\begin{proof}
The existing register $T_{\rm coef}=1+\theta\lambda$ gives
the packed mask in the form
\[
 T+1=q^2T_{\rm coef}-\ell b+\theta\ell q^4.
\]
The 94-operation schedule evaluates this expression using four
multiplications and two additions or subtractions:
\[
 \begin{aligned}
 T_{q^4}&=T_{\rm coef}q^2,&
 L_b&=\ell b,& T_2&=T_{q^4}-L_b,\\
 L_{q^4}&=\ell q^4,&
 M&=\theta L_{q^4},& T+1&=T_2+M.
 \end{aligned}
\]
Replace these six instructions by the following five:
\begin{equation}\label{eq:factored-mask-five}
 \begin{aligned}
 T_{q^4}&=T_{\rm coef}q^2,\qquad
 Q_\theta=\theta q^4,\qquad G=Q_\theta-b,\\
 I&=\ell G,\qquad T+1=T_{q^4}+I.
 \end{aligned}
\end{equation}
These use three multiplications and two additions or subtractions.
Distributivity gives the unconditional integer identity
\begin{equation}\label{eq:factored-mask-identity}
 q^2T_{\rm coef}+\ell(\theta q^4-b)
   =q^2T_{\rm coef}-\ell b+\theta\ell q^4.
\end{equation}
Thus the retained packed output, named \texttt{T} in the
certificate, has its old value for every integer assignment,
before any source equation is assumed. The deleted temporary
registers have no consumers outside this block and occur in
none of the free equality tests.

Every other instruction and equality test is retained. The
source residuals are therefore exactly those of the 94-operation
system. The identity map on the input parameters and all 34
positive supplied unknowns is the positive-witness map in both
directions. Each schedule determines its own temporary registers
from that same tuple. Consequently the full universality,
canonical-code, carry, unit, and Pell proof transfers without
any changed hypothesis.

The new temporary registers introduce no supplied unknown or
domain condition. They are in fact positive on the existing
domain: $\theta=H+b>b$, $q\ge1$, and hence
$\theta q^4-b>0$. Each subtraction is certified by its
reversed addition, as before. The complete count is
\[
 94-1=93=50\text{ multiplications}+43\text{ additions}.
\]
No numeral is introduced or changed.
\end{proof}

Appendix~\ref{app:93} gives all 93 instructions and the 22
unchanged free equality tests. The complete checker is in
\file{../verification/}:
\begin{center}
\small\file{round34\_1980\_factored\_mask\_certificate.py}.
\end{center}
It verifies the exact six-to-five replacement, the absence of
outside uses for deleted temporaries, equality of every retained
register, all primitive checks, and all source residuals with
their unchanged acyclic corrections. The shorter schedule is
a certificate for the same proved system; 93 is an upper bound,
with no optimality assertion.
